\subsubsection{3. Finite-flatness mean-field scaling}\label{finite-flatness-mean-field-scaling}

Suppose \(\mu_n\) depends on \(x\) only through \(M\), with magnetization law governed by \(p_n(m)\propto e^{-nJ(m)}\) on spacing \(1/n\). Assume \(J\) has a unique global minimum \(m_*\), and, in a neighbourhood,
\[J(m)=J(m_*)+c|m-m_*|^{2k}+o(|m-m_*|^{2k}),\qquad c>0,\]
with the usual nondegenerate Laplace-tail hypotheses. Set \(u=n^{1/(2k)}(m-m_*)\). Then
\[e^{-nJ(m)}\asymp e^{-c|u|^{2k}},\qquad dm=n^{-1/(2k)}du.\]
Consequently \(|m-m_*|=\Theta(n^{-1/(2k)})\) in probability and
\[\operatorname{Var}(M)=n^2\operatorname{Var}(m)=\Theta(n^{2-1/k}).\]
In particular \(\operatorname{Var}(M)=o(n^2)\) for every finite \(k\). The same slice-stability and small-\(n\) assembly as in Theorem 2.3 therefore proves AA for exchangeable mean-field models satisfying these regularity assumptions. The cases \(k=1,2,3\) give exponents \(1,3/2,5/3\), respectively. Multiple global minima can produce \(\operatorname{Var}(M)=\Theta(n^2)\), in which case the stability error no longer vanishes; this does not by itself prove failure for every such model.
